\section{Triple DES / TDEA kao tranziciono rešenje}
\label{sec:tdea}
% TODO: Razraditi poglavlje tek nakon provere izvora; ne unositi neproverene činjenice.
Triple DES biće razmotren kao način produženja upotrebljivosti postojeće
DES infrastrukture i kao prelaz ka novom standardu. Specifikacija i kasnija
odluka o povlačenju čine odvojene vrste dokaza~\cite{sp800-67r2,tdea-withdrawal-2023}.

\subsection{Motivacija i ograničenje Double DES-a}
% TODO: Kratko objasniti meet-in-the-middle za 2DES; ne tretirati ga kao četvrti glavni standard.

\subsection{EDE konstrukcija i kompatibilnost}
% TODO: Definisati E_K3(D_K2(E_K1(P))), DES-EDE i slučaj jednakih ključeva. Razlikovati dve nezavisne komponente (K3=K1) od tri.

\subsection{Veličina ključa i efektivna sigurnost}
% TODO: Razdvojiti 112/168 efektivnih bitova materijala, zapis sa paritetom i sigurnost pod konkretnim napadom. Dokumentovati vremenske, memorijske i podatkovne uslove; ne sabirati naivno sigurnost DES instanci.

\subsection{Performanse i prednosti tranzicije}
% TODO: Obraditi tri DES operacije, kompatibilnost, postojeći hardver i trošak migracije.

\subsection{Ograničenja bloka i Sweet32}
% TODO: Istražiti birthday granice za 64-bitni blok i konkretne uslove Sweet32 napada; ne poistovećivati koliziju bloka sa pronalaženjem ključa.

\subsection{Povlačenje i preostale nasleđene primene}
% TODO: Razdvojiti NIST zabranu primene nove zaštite od dozvoljene obrade ranije zaštićenih podataka, uz datume i opseg politike. Status standarda nije dokaz nestanka svih instalacija.
